

\section{Affine Processes on the set $S_{d}^{++}$ of strictly positive definite symmetric matrices}
\subsection{General results and notations}
To outline the setup we will consider affine processes taking values in the interior of the cone $S_d^+$. We will use the notations $\psi_t(u)=\phi(t,u)$ and $\phi_t(u)=\phi(t,u)$ so as to be consistent with \cite{article_KRPT2009}. We will be employing a property of the functions defining the Laplace transform, that we report after the following
 
%Affine processes have found many applications in Finance. For example, Gaussian interest rate models, like the ones of \cite{vas77} and \cite{hul93},  as well as square root models like in \cite{cox85}) belong to the class of affine models. \cite{heston}  developed a stochastic volatility model that can also be nested within the framewok of \cite{dar96}, who introduced in finance the notion of affine processes and allowed for both Gaussian and square root processes. In the last decade, there has been a number of theoretical and empirical applications in both interest rates and stochastic volatility modeling that systematically employed the affine structure for the stochastic factors.\\


%It is well known that it is possible to extend the notion of affine process, introduced by \cite{dar96} on the state space domain $\mathbb{R}_{\geq 0}^{m} \times \mathbb{R}^{n}$, to the domain $S_{d}^{+}$ (the set of symmetric positive semidefinite $d \times d$ matrices endowed with the scalar product $\left\langle x,y\right\rangle=Tr\left[ xy\right]$):
\begin{definition}
(\cite{article_Cuchiero}, Definition 2.1) Let $(\Omega, \mathcal{F} ,\left(\mathcal{F}_t \right)_{t \geq 0}, \mathbb{P})$ be a filtered probability space with the filtration $\left(\mathcal{F}_t \right)_{t \geq 0}$ satisfying the usual assumptions. A Markov process $\Sigma=\left(\Sigma_{t} \right)_{t \geq 0}$ with state space $S_{d}^{+}$, transition probability $p_{t}(\Sigma_0,A)=\mathbb{P}(\Sigma_{t} \in A)$ for $A \in S_{d}^{+}$, and transition semigroup $\left( P_{t}\right)_{t \geq 0}$ acting on bounded functions $f$ on $S_{d}^{+}$ is called affine process if:
\begin{enumerate}
\item it is stochastically continuous, that is, $\lim_{s \to t}p_{s}(\Sigma_0, \cdot)=p_{t}(\Sigma_0, \cdot)$ weakly on $S_{d}^{+}$ $\forall t \text{, } x \in S_{d}^{+}$, and
\item its Laplace transform has exponential-affine dependence on the initial state:
\begin{equation}
P_{t}e^{-\tr\left[u\Sigma_0 \right]}=\mathbb{E}\left[\left.e^{-\tr\left[u\Sigma_t\right]}\right|\mathcal{F}_0\right]=\int_{S_{d}^{+}}{e^{-\tr\left[u\xi \right]}p_t(\Sigma_0,d\xi)}=e^{-\phi_t(u)-\tr\left[ \psi_t(u)\Sigma_0\right]},
\label{LaplaceAffine}\end{equation}
$\forall t\text{ and } \Sigma_{0}, u \in S_{d}^{+}$, for some function $\phi: \mathbb{R}_{\geq 0}\times S_{d}^{+} \rightarrow \mathbb{R}_{\geq 0}$ and $\psi: \mathbb{R}_{\geq 0}\times S_{d}^{+} \rightarrow S_{d}^{+}$.
\end{enumerate}
\label{definition1}\end{definition}
%
%
%Note that in the definition above we assumed that the process is stochastically continuous, a feature that implies, according to Proposition 3.4 in \cite{article_Cuchiero}, that the process is regular in the following sense:
%\begin{definition} (\cite{article_Cuchiero}, Definition 2.2)
%The affine process $\Sigma$ is called regular if the derivatives
%\begin{equation}
%F(u)=\frac{\partial \phi(t,u)}{\partial t}|_{t=0^+}, \hspace{10mm}R(u)=\frac{\partial \psi(t,u)}{\partial t}|_{t=0^+}
%\end{equation}
%exist and are continuous at $u=0$.
%\label{definition2}\end{definition}
%
%
%

Having applications in mind, we will consider affine processes which are \textit{solvable} in the sense of \cite{article_Grasselli} (who investigated  affine processes on the more general symmetric cone state space domain): this means that the state space that we will consider is the interior of $S_{d}^{+}$, namely the cone of \textit{strictly positive definite} symmetric matrices, denoted by $S_{d}^{++}$\footnote{By analogy, the set of negative (resp. strictly negative) definite symmetric $d\times d$ matrices will be denoted by $S_{d}^{-}$ (resp. $S_{d}^{--}$).}.
 Solvability is important, in fact it ensures that the Riccati Ordinary Differential Equation associated to the Laplace transform (\ref{LaplaceAffine}) through the usual Feynman-Kac argument has a regular globally integrable flow: this will be crucial to outline our methodology (see e.g. the proof of Theorem \ref{TeoremaMartingale} in the sequel). \\
%
%
%  
%\cite{article_DFS} and \cite{thesis_KR} provided a full characterization of affine processes with state space $\mathbb{R}_{\geq 0}^{m} \times \mathbb{R}^{n}$ by classifying all admissibile parameters\footnote{Admissibility constitutes a weaker notion with respect to solvability since the latter is only sufficient for the existence of the process. In particular, admissibility requires technical conditions in order to understand the behavior of the process at the boundary of the state space.}. In \cite{article_Cuchiero} an analogous characterization in terms of admissibility conditions is extended to the state space $S_{d}^{+}$, thus leading to the following 
%\begin{definition}
%(\cite{article_Cuchiero} definition 2.3) An admissible parameter set $\left(\alpha,b,\beta^{ij},c,\gamma,m,\mu \right)$ associated with $\chi$\footnote{$\chi : S_d\rightarrow S_d$ denotes some bounded continuous truncation function with $\chi (\xi)=\xi$ in a neighborhood of 0.} consists in: a linear diffusion coefficient $\alpha \in S_{d}^{+}$, a constant drift term $b\succeq (d-1)\alpha$\footnote{For $a,b\in S_{d}^{+}$ we write $a\succeq b$ if $a-b\in S_{d}^{+}$.}, a constant killing rate term $c \in \mathbb{R}^{+}$, a linear killing rate coefficient $\gamma \in S_{d}^{+}$.
%Moreover we admit jumps satisfying the following conditions:
%\begin{itemize}
%\item a constant jump term: a Borel measure $m$ on $S_{d}^{+} \setminus \left\{ 0\right\}$ satisfying
%\begin{equation}
%\int_{S_{d}^{+} \setminus \left\{ 0\right\}}{\left(\parallel \xi \parallel \land 1 \right)m(d \xi)}<\infty
%\end{equation}
%\item a linear jump coefficient: a $d \times d$ matrix $\mu=(\mu_{ij})$ of finite signed measure on $S_{d}^{+} \setminus \left\{ 0\right\}$ such that $\mu(E)\in S_{d}^{+}$ $\forall E \in \mathcal{B}(S_{d}^{+}\setminus\left\{0\right\})$ and the kernel
%\begin{equation}
%M(x,d\xi):=\frac{Tr\left[x \mu(d\xi) \right]}{\parallel \xi \parallel^2 \land 1}
%\end{equation}
%satisfies
%\begin{equation}
%\int_{S_{d}^{+} \setminus \left\{ 0\right\}}{Tr\left[\chi(\xi)u \right]M(x,d\xi)}<\infty \hspace{10mm} \forall x,u \in S_{d}^{+}\text{ s.t. }Tr\left[xu\right]=0
%\end{equation}
%\end{itemize}
%Finally we have a linear drift coefficient: a family $\beta^{ij}=\beta^{ji} \in S_{d}$ s.t. the linear map $B:S_{d} \rightarrow S_{d}$ of the form 
%\begin{equation}
%B(x)=\sum_{i,j}{\beta^{ij}x_{ij}}
%\end{equation}
%satisfies
%\begin{equation}
%Tr\left[B(x)u\right]-\int_{S_{d}^{+} \setminus \left\{ 0\right\}}{Tr\left[\chi(\xi)u \right]M(x,d\xi)}\geq 0 \hspace{5mm}\forall x,u \in S_{d}^{+} \text{ with } Tr\left[xu\right]=0.
%\end{equation}
%\end{definition}
%The following theorem is a generalization of the result by \cite{article_DFS} to the state space $S_{d}^{+}$. Denote by $\mathcal{S}_{+}$ the space of rapidly decreasing $C^\infty$-functions on $S_d^+$ and let $\mathcal{D}(\mathcal{A})$ be the domain of the generator of the process.
%
%\begin{theorem}
%(\cite{article_Cuchiero} Theorem 2.4) Suppose $\Sigma$ is an affine process on $S_{d}^{+}$. Then $\Sigma$ is regular and has the Feller property. Let $\mathcal{A}$ be its infinitesimal generator on $C_{0}(S_{d}^{+})$. Then $\mathcal{S}_{+} \subset \mathcal{D}(\mathcal{A})$ and there exists an admissible parameter set $\left(\alpha,b,\beta^{ij},c,\gamma,m,\mu \right)$ such that, for $f \in \mathcal{S}_{+}$
%\begin{align}
%\mathcal{A}f(x)&=\frac{1}{2}\sum_{i,j,k,l}{A_{ijkl}(x)\frac{\partial^{2}f(x)}{\partial x_{ij} \partial x_{kl}}}+\sum_{i,j}\left(b_{ij}+\beta_{ij}(x) \right)\frac{\partial f(x)}{\partial x_{ij}}-\left(c+Tr\left[ \gamma x\right] \right)f(x)\nonumber\\
%&+\int_{S_{d}^{+} \setminus \left\{ 0\right\}}{\left(f(x+\xi)-f(x) \right)m(d\xi)}\nonumber\\
%&+\int_{S_{d}^{+} \setminus \left\{ 0\right\}}{\left(f(x+\xi)-f(x)-Tr\left[\chi(\xi)\nabla f(x)\right] \right)M(x,d\xi)},\label{generatoraffine}
%\end{align}
%where 
%\begin{equation}
%A_{ijkl}(x)=x_{ik}\alpha_{jl}+x_{il}\alpha_{jk}+x_{jk}\alpha_{il}+x_{jl}\alpha_{ik}.
%\end{equation}
%Moreover, $\phi(t,u)$ and $\psi(t,u)$ in definition (\ref{definition1}) solve the generalized Riccati differential equations, for $u \in S_{d}^{+}$,
%\begin{align}
%\frac{\partial \phi(t,u)}{\partial t}&=F(\psi(t,u)), \hspace{10mm} \phi(0,u)=0,\nonumber\\
%\frac{\partial \psi(t,u)}{\partial t}&=R(\psi(t,u)), \hspace{10mm} \phi(0,u)=u,\nonumber
%\label{ODE}
%\end{align}
%with
%\begin{align}
%F(u)&=Tr\left[bu\right]+c-\int_{S_{d}^{+} \setminus \left\{ 0\right\}}{\left( e^{-Tr\left[u\xi\right]}-1\right)m(d\xi)},\nonumber\\
%R(u)&=-2u\alpha u+ \beta^{T}(u)+\gamma\nonumber\\
%&-\int_{S_{d}^{+} \setminus \left\{ 0\right\}}{\left(\frac{e^{-Tr\left[u\xi\right]}-1-Tr\left[\chi(\xi)u\right]}{\parallel \xi \parallel^2 \land 1}\right)\mu(d\xi)},
%\end{align}
%where $B^{T}_{ij}(u)=Tr\left[\beta^{ij}u\right]$.\\
%Conversely, let $\left(\alpha,b,\beta^{ij},c,\gamma,m,\mu \right)$ be an admissible parameter set. Then there exists a unique affine process on $S_{d}^{+}$ with infinitesimal generator given by (\ref{generatoraffine}) and such that the affine transform formula (\ref{LaplaceAffine}) holds for all $(t,u)\in \mathbb{R}_{\geq 0} \times S_{d}^{+}$ where $\phi(t,u)$ and $\psi(t,u)$ are given by (\ref{ODE}).
%\end{theorem}




The next property closes our survey on affine processes. It will be needed when we prove that the structure of the model is preserved under changes of measure.
\begin{lemma}(\cite{article_Cuchiero} Lemma 3.2) Let $\Sigma$ be an affine process on $S_d^+$, then the functions $\phi$ and $\psi$ satisfy the following property:
\begin{align}
\phi_{t+s}(u)&=\phi_t(u)+\phi_s(\psi_t(u)),\nonumber\\
\psi_{t+s}(u)&=\psi_s(\psi_t(u)).\nonumber
\end{align}
\end{lemma}

\subsection{Examples}

The previous general framework may be quite abstract at a first sight, mostly because of the high technical level required to properly introduce the notion of admissibility and existence for affine processes, see \cite{article_Cuchiero}. In this subsection we provide some examples which will illustrate some concrete applications. We start with the most important one, which will also constitute our main object of study in the numerical illustrations.

\subsubsection{The Wishart process}\label{Wis_disc}

We suppose that the process $\Sigma$ is governed by the following (matrix) SDE:
\begin{equation}
d\Sigma_t=(\Omega\Omega^{\top}+M\Sigma_t+\Sigma_tM^{\top})dt+\sqrt{\Sigma_t}dW_tQ+Q^{\top}dW_t^{\top}\sqrt{\Sigma_t},\label{wishart}
\end{equation}
which was first studied by \cite{article_Bru} and whose solution is known as Wishart process. We assume $M,Q$ invertible and $M$ negative definite so as to ensure stationarity of the process. Moreover we require $\Omega\Omega^{\top}=\kappa Q^{\top}Q$ for a real parameter $\kappa \geq d+1$ to grant solvability (or equivalently to grant that $Det(\Sigma_t)>0$ with probability 1). Under the solvability assumption \cite{article_Grasselli} showed that the Riccati ODE corresponding to the characteristic function can be linearized and therefore admits a closed form solution. This is important in view of possible applications since in this case the functions $\phi$ and $\psi$ in definition (\ref{definition1}) are explicitly known:
\begin{proposition}
Consider the process $\Sigma=\left(\Sigma_t\right)_{0\leq t\leq T}$ which solves the SDE \eqref{wishart}. Then the conditional Laplace transform is given by:
\begin{align}
\mathbb{E}\left[\left.e^{-\tr\left[u\Sigma_T\right]}\right|\mathcal{F}_t\right]=e^{-\phi_\tau(u)-\tr\left[\psi_\tau(u)\Sigma_t\right]},
\end{align}
where $\tau:=T-t$. The functions $\phi_\tau(u)$ and $\psi_\tau(u)$ satisfy the following system of ODE's:
\begin{align}
\frac{\partial \psi_{\tau}}{\partial \tau}&=\psi_\tau(u) M+M^\top \psi_\tau(u)-2\psi_\tau(u) Q^\top Q\psi_\tau(u),\quad\psi_0(u)=u,\label{Riccati_psi}\\
\frac{\partial \phi_{\tau}}{\partial \tau}&=\tr\left[\kappa Q^{\top}Q\psi_\tau(u)\right],\quad\phi_0(u)=0\label{Riccati_phi}
\end{align}
which is solved by
\begin{align}
\psi_{\tau}(u)=\label{clever2}\left(u\psi_{12,\tau}(u)+\psi_{22,\tau}(u) \right)^{-1}\left(u\psi_{11,\tau}(u)+\psi_{21,\tau}(u)\right), 
\end{align}
where
\begin{equation}
\left(\begin{array}{rr}
\psi_{11,\tau}(u) & \psi_{12,\tau}(u)\\
\psi_{21,\tau}(u) & \psi_{22,\tau}(u)
\end{array}\right)
=\exp\left\{ \tau \left( 
\begin{array}{rr}
M & 2Q^{\top}Q\\
0 & -M^{\top}
\end{array}
\right) \right\}
\end{equation} 
and
\begin{equation}
\phi_{\tau}(u)=\frac{\kappa}{2}\tr\left[\log\left(u\psi_{12,\tau}(u)+\psi_{22,\tau}(u) \right)+M^{\top}\tau \right].
\end{equation}
\label{proplinearization}\end{proposition}
\begin{proof}
See  \cite{article_Grasselli}.
\end{proof}
The Wishart process constitutes the matrix analogue of the square root (Bessel) process. In fact we have that the matrix $M$ can be thought of as a mean reversion parameter: this is evident from the Lyapunov equation defining the long-run matrix $\Sigma_\infty$, which is given by
\begin{align}
-\kappa Q^{\top}Q=M\Sigma_\infty+\Sigma_\infty M^{\top}.\label{Lyap}
\end{align}

The second way to appreciate the analogies w.r.t the square root process is to look at the dynamics of the entries of the matrix process $\Sigma$. Concentrating on the main diagonal, in the $2\times 2$ case we have:
\begin{align}
d\Sigma_{11}&=\left(\kappa\left(Q_{11}^2+Q_{21}^2\right)+2\left(M_{11}\Sigma_{11}+M_{12}\Sigma_{12}\right)\right)dt\nonumber\\
&+2\sigma_t^{11}\left(Q_{11}dW_t^{11}+Q_{21}dW_t^{12}\right)+2\sigma_t^{12}\left(Q_{11}dW_{21}+Q_{21}dW_t^{22}\right)\\
d\Sigma_{22}&=\left(\kappa\left(Q_{22}^2+Q_{12}^2\right)+2\left(M_{21}\Sigma_{12}+M_{22}\Sigma_{22}\right)\right)dt\nonumber\\
&+2\sigma_t^{12}\left(Q_{12}dW_t^{11}+Q_{22}dW_t^{12}\right)+2\sigma_t^{22}\left(Q_{12}dW^{21}+Q_{22}dW^{22}\right)
\end{align}
where we set
\begin{align}
\left(\begin{array}{rr}
\sigma_{11}&\sigma_{12}\\
\sigma_{12}&\sigma_{22}
\end{array}\right):=\sqrt{\Sigma}.
\end{align}

We refer to \cite{article_DaFonseca4} for additional insights on the behavior of the Wishart process when aggregating its parameters.




\subsubsection{The pure jump OU process}\label{pure_jump}
The procedure we adopt in this paper is general, meaning that we can consider different examples of processes lying in the cone of positive definite matrices. In particular, we may consider the matrix subordinators proposed by \cite{article_BNS02}, or jump-diffusions like in \cite{article_LT}. In what follows we provide some examples with the calculations of the functions $\phi_{\tau}$ and $\psi_{\tau}$.\\

Let us consider the SDE
\begin{align}
d\Sigma_t=M\Sigma_t+\Sigma_tM^{\top}+dL_t,
\end{align}
where $M\in GL(d)$ is assumed as usual to be negative definite so as to grant stationarity, and $L_t$ is a pure jump process (compound Poisson Process) with constant intensity $\lambda$ and jump distribution $\nu$ with support on $S_d^{++}$.
The strong solution to this equation is given by:
\begin{align}
\Sigma_t=e^{Mt}\Sigma_0e^{M^{\top}t}+\int_{0}^{t}{e^{M(t-s)}dL_se^{M^{\top}(t-s)}}.
\end{align}
We are interested in the computation of the Laplace transform of this family of processes:
\begin{align}
\mathbb{E}\left[\left.e^{\tr\left[u\Sigma_T\right]}\right|\mathcal{F}_t\right]=e^{-\phi_\tau(u)-\tr\left[\psi_\tau(u)\Sigma_t\right]}.
\end{align}
The functions $\phi$ and $\psi$ solve the following (matrix) ODE's:
\begin{align}
\frac{\partial \psi_{{\tau}}}{\partial \tau}&=\psi_\tau(u)M+M^{\top}\psi_\tau(u)\quad\psi_0(u)=u\\
\frac{\partial \phi_{{\tau}}}{\partial \tau}&=-\lambda\int_{S_d^+\setminus\left\{0\right\}}{\left(e^{-\tr\left[\psi_\tau(u)\xi\right]}-1\right)\nu(d\xi)}\quad\phi_0(u)=0.
\end{align}
The solution for the first ODE is given by:
\begin{align}
\psi_\tau(u)=e^{M^{\top}\tau}ue^{M\tau},
\end{align}
so we can compute the Laplace transform by quadrature:
\begin{align}
\frac{\partial \phi_{\tau}}{\partial \tau}&=-\lambda\int_{S_d^+\setminus\left\{0\right\}}{\left(e^{-\tr\left[e^{M^{\top}\tau}ue^{M\tau}\xi\right]}-1\right)\nu(d\xi)}.
\end{align}
In the following we provide explicit computations by assuming some particular distribution $\nu(\cdot)$ for the jump size. The proofs of this formulae may be found in \cite{book_GuptaNagar2000}. For the sake of clarity, we specify that the Wishart distribution that we consider in the next sections are the classical distributions arising in the context of multivariate statistics.\\

\paragraph{\textbf{Wishart Distribution.}}
Let $J$ be the jump size. Consider the case $J\sim Wis_d\left(n,\mathcal{Q}\right)$. Then we have
\begin{align}
\phi_\tau(u)=-\lambda\int_{0}^{\tau}{\det\left(I_d+2e^{M^{\top}s}ue^{Ms}\mathcal{Q}\right)^{-\frac{n}{2}}ds}+\lambda \tau.
\end{align}
\paragraph{\textbf{Non-Central Wishart Distribution.}}
Let be $J\sim Wis_d\left(n,\mathcal{Q},\mathcal{M}\right)$, then we have
\begin{align}
\phi_\tau(u)&=-\lambda\int_{0}^{\tau}{\det\left(\mathcal{Q}\right)^{-\frac{n}{2}}\det\left(2e^{M^{\top}s}ue^{Ms}+\mathcal{Q}^{-1}\right)^{-\frac{n}{2}}}\times\nonumber\\
&\exp\left\{\tr\left[-\frac{1}{2}\mathcal{Q}^{-1}\mathcal{M}\mathcal{M}^{\top}+\frac{1}{2}\mathcal{Q}^{-1}\mathcal{M}\mathcal{M}^{\top}\mathcal{Q}^{-1}\left(2e^{M^{\top}s}ue^{Ms}+\mathcal{Q}^{-1}\right)\right]\right\}ds\nonumber\\
&+\lambda \tau.
\end{align}
\paragraph{\textbf{Beta type I distribution.}} Let be $J\sim\beta_d^I(a,b)$, then
\begin{align}
\phi_\tau(u)&=-\lambda\int_{0}^{\tau}{_{1}F_1(a;a+b;-e^{M^{\top}s}ue^{Ms})ds}+\lambda \tau.
\end{align}

\paragraph{\textbf{Beta  type II distribution.}} Let be $J\sim\beta_d^{II}(a,b)$, then
\begin{align}
\phi_\tau(u)&=-\lambda\int_{0}^{\tau}{\frac{\Gamma_d(a+b)}{\Gamma_d(b)}\Psi(a;-b+\frac{1}{2}(d+1);e^{M^{\top}s}ue^{Ms})ds}+\lambda \tau,
\end{align}
where $_{m}F_{n}$, $\Gamma_d(a)$, and $\Psi(a;b;R)$ denote respectively the hypergeometric function of matrix argument, the multivariate Gamma function and the confluent hypergeometric function, see e.g. \cite{book_GuptaNagar2000}.


